% Generated from locale/tr/content/lambda-calculus/church-rosser/parallel-beta-eta-reduction.tex; do not hand-edit.


\olfileid[tr]{lam}{cr}{pbe}

\olsection{Paralel $\beta\eta$-indirgeme}

Bu bölümde $\beta\eta$-indirgemeye karşılık gelen paralel indirgeme kavramı
olan paralel $\beta\eta$-indirgemenin Church--Rosser özelliğini ispatlıyoruz.

\begin{defn}[Paralel $\beta\eta$-indirgeme, $\beredpar$]
  \ollabel{defn:beredpar}
  Terimler üzerindeki \emph{paralel $\beta\eta$-indirgeme} ($\beredpar$)
  tümevarımlı olarak şöyle tanımlanır:
  \begin{enumerate}
    \item \ollabel{defn:beredpar1} $x \beredpar x$.
    \item \ollabel{defn:beredpar2} $N \beredpar N'$ ise
      $\lambd[x][N] \beredpar \lambd[x][N']$.
    \item \ollabel{defn:beredpar3} $P \beredpar P'$ ve $Q \beredpar Q'$ ise
      $PQ \beredpar P'Q'$.
    \item \ollabel{defn:beredpar4} $N \beredpar N'$ ve $Q \beredpar Q'$ ise
      $(\lambd[x][N])Q \beredpar \Subst{N'}{Q'}{x}$.
    \item \ollabel{defn:beredpar5} $N \beredpar N'$ ise ve
      $x \notin \FV{N}$ olmak koşuluyla
      $\lambd[x][Nx] \beredpar N'$.
  \end{enumerate}
\end{defn}

\begin{thm}\ollabel{thm:refl}
  $M \beredpar M$.
\end{thm}

\begin{proof}
  Alıştırma.
\end{proof}

\begin{prob}
  \olref[lam][cr][pbe]{thm:refl} teoremini ispatlayın.
\end{prob}

\begin{defn}[$\beta\eta$-tam geliştirme]\ollabel{defn:becd}
  $M$ teriminin \emph{$\beta\eta$-tam geliştirmesi} $\becd{M}$ şöyle
  tanımlanır:
  \begin{align}
    \becd{x} &= x \ollabel{defn:becd1} \\
    \becd{(\lambd[x][N])} &= \lambd[x][\becd{N}] \ollabel{defn:becd2}\\
    \becd{(PQ)} &= \becd{P}\becd{Q} && \text{$P$ bir $\lambd$-soyutlaması değilse}
    \ollabel{defn:becd3} \\
    \becd{((\lambd[x][N])Q)} &= \Subst{\becd{N}}{\becd{Q}}{x}
                             \ollabel{defn:becd4} \\
    \becd{(\lambd[x][Nx])} &= \becd{N} \ollabel{defn:becd5} &&
                             \text{$x\notin\FV{N}$ ise.}
  \end{align}
\end{defn}

\begin{lem}\ollabel{lem:comp}
  $M \beredpar M'$ ve $R \beredpar R'$ ise $\Subst{M}{R}{y}
  \beredpar \Subst{M'}{R'}{y}$.
\end{lem}

\begin{proof}
  $M \beredpar M'$ !!{derivation} üzerinde tümevarımla.

  İlk dört durum \olref[pb]{lem:comp} içindeki durumların aynısıdır. Son kural
  \olref{defn:beredpar5} ise bazı $x,N,N'$ için
  $M=\lambd[x][Nx]$, $M'=N'$, $x\notin\FV{N}$ ve $N\beredpar N'$ olur.
  Göstermek istediğimiz
  $\Subst{(\lambd[x][Nx])}{R}{y}\beredpar\Subst{N'}{R'}{y}$, yani
  $\lambd[x][\Subst{N}{R}{y}x]\beredpar\Subst{N'}{R'}{y}$ bağıntısıdır.
  Bu, \olref{defn:beredpar}\olref{defn:beredpar5} ile tümevarım
  varsayımından çıkar.
\end{proof}

\begin{lem}\ollabel{lem:cont}
  $M \beredpar M'$ ise $M' \beredpar \becd{M}$.
\end{lem}

\begin{proof}
  $M \beredpar M'$ !!{derivation} üzerinde tümevarımla.

  İlk dört durum \olref[pb]{lem:cont} içindeki durumlara benzer. Son kural
  \olref{defn:beredpar5} ise bazı $x,N,N'$ için
  $M=\lambd[x][Nx]$, $M'=N'$, $x\notin\FV{N}$ ve $N\beredpar N'$ olur.
  Göstermek istediğimiz $N'\beredpar\becd{(\lambd[x][Nx])}$, yani
  $N'\beredpar\becd{N}$ bağıntısıdır; bu da tümevarım varsayımından doğrudan
  çıkar.
\end{proof}

\begin{thm}\ollabel{thm:cr}
  $\beredpar$ Church--Rosser özelliğini sağlar.
\end{thm}

\begin{proof}
  Doğrudan \olref{lem:cont} lemasından çıkar.
\end{proof}

